%% ============================================================
%%  Appendix D — Chapter 5 (Temporal Computing) Supporting Info
%%
%%  \include'd from thesis.tex under \appendix. Documents the
%%  parameter-card sourcing and the simulation / read-out
%%  settings, so the in-silico results are reproducible from the
%%  measured archive without consulting the scripts.
%% ============================================================

\chapter{Supporting Information for Chapter~\ref*{ch:temporal}}\label{app:ch5_si}

This appendix documents how the behavioural-model parameter cards of
\cref{sec:ch5_model} are sourced from the measured archive, and lists the
simulation and read-out settings used across the three evidence tiers. Together
with the per-cell tables of \cref{app:ch4_celltables}, it makes the in-silico
results reproducible without consulting the analysis scripts
(\path{scripts/ch5_figures.py}, \path{scripts/ch5_physio_context.py},
\path{scripts/ch5_onset.py}).

\section{Parameter-Card Sourcing}\label{app:ch5_cards}

Each composition cell is reduced to a parameter card
$(\tau_c,\beta_c,\alpha_c,\text{peak}_c)$ that drives one reservoir node type. The
values are taken directly from the comparative fits, with no free re-tuning:
\begin{itemize}
  \item the fading-memory law $\lambda_c(\Delta t)=\exp[-(\Delta t/\tau_c)^{\beta_c}]$ uses the per-cell $(\tau_c,\beta_c)$ of \cref{tab:ch4_decay_by_cell}; for the few cells without an identified stretched fit, a single exponential reproducing the model-free half-life is used instead ($\tau_c=t_{1/2}/\ln 2$, $\beta_c=1$);
  \item the write nonlinearity uses the per-cell growth exponent $\alpha_c$ and peak enhancement of \cref{tab:ch4_pulses_by_cell};
  \item the device-to-device spread injected into each bank is the within-cell scatter of those same tables.
\end{itemize}
The heterogeneous banks of \cref{sec:ch5_benchmarks,sec:ch5_physio} span the
replicated \PEO{}/\LiTf{} composition cells of \cref{tab:ch4_decay_by_cell}
($\tau$ from $\approx 3$ to $\approx 26$\,s); the homogeneous control banks are
jittered copies of the lead cell \PEO{}~0.3/0.09 alone.

\section{Simulation and Read-out Settings}\label{app:ch5_settings}

\Cref{tab:ch5_settings} collects the settings used throughout
\cref{ch:temporal}. All read-outs are a single linear ridge regression; no
nonlinear decoder is used, so any demonstrated computation is attributable to the
device nodes rather than the read-out.

\begin{table}[t]
  \centering
  \caption[Reservoir simulation and read-out settings]{Simulation and read-out settings across the three evidence tiers of \cref{ch:temporal}. All banks are spatial banks of independent (non-recurrent) leaky nodes driven by randomly-masked input; all read-outs are linear ridge regression.}
  \label{tab:ch5_settings}
  \footnotesize
  \setlength{\tabcolsep}{5pt}
  \begin{tabularx}{\linewidth}{@{}>{\raggedright\arraybackslash}p{3.4cm}X@{}}
    \toprule
    Setting & Value \\
    \midrule
    Node model & $x^{(i)}_n=\lambda_i x^{(i)}_{n-1}+[\max(w_i u_n,0)]^{\alpha_i}$ (\cref{sec:ch5_model}) \\
    Read-out & single linear ridge regression \\
    Bank size $N$ & 24 (benchmarks); 48 (physiology/affect) \\
    Seeds & 10 random input/bank seeds; means $\pm$ SD reported \\
    Significance tests & paired Wilcoxon signed-rank over seeds (or subjects) \\
    Memory capacity & $\mathrm{MC}_k$ summed over lags, random uniform input \cite{Jaeger2002STM} \\
    NARMA & NARMA-10 system identification, NRMSE \cite{AtiyaParlos2000} \\
    Information-processing capacity & degree-1 (linear) + degree-2 (nonlinear), bounded by $N$ \cite{Dambre2012} \\
    Affective corpus & \WESAD{} chest signals, leave-one-subject-out over 15 subjects \cite{Schmidt2018WESAD} \\
    Streaming cadence & $\Delta t=\SI{1}{\second}$; reconstruction delays $1,3,8,20,45$\,s \\
    Affect scoring & class-balanced macro-F1 \\
    \bottomrule
  \end{tabularx}
\end{table}

\section{Continuous-Monitoring Stress Test}\label{app:ch5_robust}

The deployment-robustness result of \cref{sec:ch5_robust} (\path{scripts/ch5_onset.py})
is a continuous binary stress detector: at every $\Delta t = \SI{1}{\second}$ step the
read-out predicts stress (\WESAD{} label~2) versus not-stress (baseline and amusement
pooled), leave-one-subject-out, with the same single linear ridge read-out and
$15$-second causal label smoothing used in the streaming task. Four matched banks of
$N=24$ nodes are compared on identical input masks: the instantaneous input (no
memory), a memoryless bank (leak set to zero, dimensionality only), the homogeneous
lead-only bank, and the heterogeneous composition bank.

Sensor corruption is modelled by adding, to the per-subject scaled streams
(channels in $[0,1]$), independent Gaussian noise of standard deviation $\sigma$ at
every sample, plus sparse motion-artefact bursts (short windows of amplified noise at
a $1\%$ per-sample rate), clipped back to $[0,1]$; the same realisation is applied to
the training and held-out streams so the comparison isolates each architecture's
intrinsic noise tolerance. Robustness curves are seed-averaged over three bank/noise
seeds; the per-subject significance test at $\sigma=0.4$ is a paired Wilcoxon
signed-rank over the fifteen subjects. Reported metrics are the binary macro-F1, the
pooled false-alarm rate (fraction of not-stress steps called stress), and --- for
completeness, though they do not separate the banks --- the transition-window
macro-F1 (steps within $\pm\SI{20}{\second}$ of a label change) and the onset
detection latency (time to the first $\geq\SI{5}{\second}$ sustained stress call after
each onset).
